\section{结论}
\label{sec:concl}

本文主张，不完整历史应被视为一种预测条件，而非有待修复的数据，并以一个协议与一个模型支撑该主张。TrajGap-Bench 在公开 TrajImpute 发布上冻结单一 Mixed 契约——一种划分构造、一条检查点规则、两项指标、五个诊断分层——使所有消费不完整历史的方法可以在相同条件下比较。GAST 使观测证据状态贯穿逐帧关系注意力、可靠度调制的因果状态空间编码器与锚点条件化目标解码器，且从不对坐标做插补。

在五个场景共 302{,}886 个测试样本上，GAST 取得对比方法中最优的宏平均 best-of-20 误差（0.3168/0.4831~m），相对最强外部基线 MoFlow 降低 9.4\%/7.6\%，并在 $\mathrm{minADE}_{20}$ 上对 MoFlow 五个场景全胜、在 $\mathrm{minFDE}_{20}$ 上四胜一负。该优势与证据条件化的设计意图一致：相对 MoFlow 的 $\mathrm{minFDE}_{20}$ 降低幅度由无缺失时的 5.4\% 增至缺失六帧时的 22.6\%，并在仅剩一个观测时反转。本文同时报告了两项限定该贡献的结果：改进并不跨场景一致（HOTEL-M 与 ZARA1-M 的 $\mathrm{minFDE}_{20}$ 分别落后于 AgentFormer 与 MoFlow），且 best-of-20 的提升未同步体现在概率最高的单一模式上。

下一步工作直接来自上述局限。需要多种子研究与按轨迹聚类的统计，以确定相对外部基线的边际是否稳定；需要同一协议下逐组件开关的消融，以分离逐帧几何证据编码、可靠度调制时间聚合与锚点条件化解码的贡献；需要为证据年龄引入饱和规则以处理最旧锚点处的反转；并需要改进模式评分目标，使覆盖度收益能够传递到 top-1。上述协议亦应扩展到其他观测几何与真实感知流程产出的轨迹，后者的缺失分布并非合成。
